% Overleaf-ready manuscript: tone-routed flash/no-flash skin colorimetry
% Upload to Overleaf (pdfLaTeX). Merge \bibliography{} with your project .bib.
% Code & data: https://github.com/RooneyEmily/Fitskin
%
% Compile: pdflatex skin_tone_calibration_overleaf.tex
%          bibtex skin_tone_calibration_overleaf   (if using external .bib)

\documentclass[11pt]{article}

\usepackage[margin=1in]{geometry}
\usepackage{booktabs}
\usepackage{siunitx}
\usepackage{amsmath}
\usepackage{graphicx}
\usepackage[hidelinks]{hyperref}
\usepackage{enumitem}
\usepackage{algorithm}
\usepackage{algpseudocode}

\sisetup{
  round-mode=places,
  round-precision=2,
  detect-weight=true,
  detect-family=true,
}

\newcommand{\deee}{\ensuremath{\Delta E_{00}}}
\newcommand{\Maff}{\mathbf{M}_{\mathrm{aff}}}
\newcommand{\Rvec}{\mathbf{R}}

\title{Skin-Tone--Routed Flash/No-Flash Colorimetry from iPhone ProRAW:\\
  Offline Calibration, In-Scene ColorChecker, and Chart-Free Deployment}
\author{%
  % TODO: author list
}
\date{June 2026}

\begin{document}
\maketitle

\begin{abstract}
Consumer iPhone photography offers a practical route to non-contact skin colorimetry, but a single camera-to-XYZ calibration matrix fails across the lightness range spanned by lighter and darker skin tones.
We evaluate flash/no-flash ($\sqrt{A\odot B}$) reflectance recovery against FitSkin contact-scanner cheek CIELAB on a two-participant iPhone ProRAW cohort (Emily, lighter; Likitha, Indian, darker skin), with cross-session ground truth.
Three pipelines are compared: (i)~in-scene Macbeth ColorChecker correction with automatic skin-tone routing of region-of-interest and matrix form; (ii)~chart-free inference using offline matrices trained on separate ColorChecker sessions; and (iii)~a hybrid production rule (chart when visible, flash/no-flash otherwise).
Fixed cheek-region extraction on chart-free darker skin inflated $\deee$ to ${\sim}9$; routing to a full face-mesh region and a tone-specific affine matrix reduced Likitha's chart-free median $\deee$ from ${\sim}5.0$ to ${\sim}2.2$ while preserving Emily's performance (${\sim}1.7$--$2.4$).
With ColorChecker visible, automatic tone routing achieved median $\deee=1.41$ ($N=6$).
We further show that training-data composition is tone-dependent: stacking booth RAW and JPEG sessions helps the lighter-skin matrix but degrades darker-skin ProRAW inference when capture pipelines differ.
An ancestry-appropriate ISSA spectral prior (\texttt{issa\_median\_south\_asian} for Likitha) marginally improves the dark bundle over an African prior.
These results support commercial shade-matching and longitudinal skincare applications at $\deee\approx 2$--$4$ without a per-capture ColorChecker, while FitSkin-grade absolute colorimetry still benefits from an in-scene reference or same-session contact pairing.
\end{abstract}

%=============================================================================
\section{Introduction}\label{sec:intro}
%=============================================================================

Skin tone measurement from camera imagery has broad applications in clinical dermatology, beauty industry color matching, and fairness auditing of face analysis systems across demographic groups.
Accurate colorimetric recovery from consumer cameras has been pursued since the early digital imaging era \cite{romero2006spectral}, yet the problem remains practically unsolved in unconstrained settings: existing pipelines typically require a physical calibration target or dedicated illuminant sensor at acquisition time, neither of which can be assumed in consumer deployments.
Classical illuminant estimation relies on statistical assumptions over scene reflectance that are routinely violated in practice \cite{gijsenij2011}.
Learning-based methods reduce this dependence but introduce their own requirements: large paired training datasets and, critically, a calibration reference at deployment, limiting their applicability outside controlled laboratory conditions.

Flash photography offers an alternative.
By capturing a matched ambient and flash image pair, the flash contribution can be isolated and used to recover surface reflectance without statistical assumptions about global scene reflectance distributions \cite{dicarlo2001illuminating,lu2006practical}.
Lu and Drew~\cite{lu2006practical} showed that ambient illuminant chromaticity can be estimated from flash/no-flash pairs in log geometric-mean chromaticity space without prior knowledge of flash spectral power distribution or camera responsivities.
Subsequent work structured flash/no-flash pairs through intrinsic decomposition~\cite{maralan2023computational} and large-scale learned decomposition on FAID~\cite{flashambient,upadhyay2025low}.
However, targetless flash/no-flash colorimetry of \emph{human skin}---evaluated against instrument-grade contact measurement across tone---has received limited systematic study.

Skin colour pipelines are additionally sensitive to how calibration is performed.
Harville et al.~\cite{harville2005consistent} and Choi et al.~\cite{choi2017performance} showed that ColorChecker matrices weighted toward skin-coloured patches outperform general 24-patch fits in the skin tone region, implying that both \emph{where} skin is sampled and \emph{which training reflectances} enter the fit matter.
Ordinal scales such as Fitzpatrick~\cite{fitzpatrick1988} underrepresent darker tones and lack device-independent coordinates~\cite{cook2025colorimetric}; the Monk Skin Tone (MST) scale~\cite{cheng2024monastic} and dedicated colorimetric skin scales~\cite{cook2025colorimetric} provide reference Lab values better suited to objective evaluation.
The ISSA dataset supplies median spectral reflectance priors stratified by ancestry~\cite{issa2025}, relevant both to learning-based auto white balance~\cite{zhou2025scr,wang2025perform} and to anchoring offline camera calibration.

This paper extends the flash/no-flash lineage to \textbf{tone-routed, FitSkin-validated skin colorimetry} on iPhone ProRAW.
We separate two decisions that prior work often conflates: (a)~\emph{skin-tone tier}---a lightness-based routing of ROI and matrix form inferred from an uncorrected preview; and (b)~\emph{ancestry spectral prior}---ISSA median reflectance used as a synthetic training row and SCR-AWB prior.
We train offline ColorChecker matrices \emph{without} requiring a chart at inference, evaluate against cross-session FitSkin cheek Lab, and report when chart-free deployment is commercially plausible.

\paragraph{Contributions.}
\begin{itemize}[noitemsep]
  \item A \textbf{skin-tone auto policy} for in-scene ColorChecker capture: preview cheek $L^*$ routes lighter skin to cheek ROI with a $3\times3$ matrix and darker skin to face-mesh ROI with a $3\times4$ affine map, improving median $\deee$ from $4.64$ to $1.41$ on six Pansor ProRAW trials.
  \item \textbf{Tone-specific offline calibration bundles} (\texttt{light/}, \texttt{dark/}) for chart-free flash/no-flash inference, reducing darker-skin chart-free median $\deee$ from ${\sim}5.0$ to ${\sim}2.2$ without regressing lighter skin.
  \item A \textbf{training-data policy} showing that multi-session stacking benefits lighter-skin calibration but must respect capture-pipeline homogeneity for darker ProRAW inference.
  \item \textbf{Ancestry-appropriate ISSA priors} (South Asian vs.\ Caucasian) distinct from lightness-tier routing, with validation on an Indian participant.
  \item A \textbf{commercial feasibility analysis} distinguishing shade matching, longitudinal skincare, and FitSkin-grade absolute colorimetry.
\end{itemize}

%=============================================================================
\section{Background}\label{sec:background}
%=============================================================================

\subsection{Camera-Based Colorimetry}

Digital cameras have been investigated as practical imaging colorimeters since the late 1990s.
Wu et al.~\cite{wu1999} demonstrated that calibration matrices mapping filtered camera RGB to CIE XYZ under selected illumination could yield spatially resolved tristimulus data.
Accurate recovery requires sensor spectral sensitivities to approximate a linear transform of the CIE colour matching functions (Luther--Ives); consumer iPhone sensors violate this, necessitating scene- or session-specific calibration~\cite{hubel1994comparison,wang2021optimized}.
Target-based ColorChecker calibration~\cite{wakholi2026systematic,BabelColorMCC} captures camera response and illumination jointly but requires the chart in frame---impractical for daily consumer capture.

\subsection{Illuminant Estimation}

Grey-world, grey-edge, white-patch, and achromatic-pixel methods encode reflectance priors that fail on portrait scenes~\cite{gijsenij2011}.
Learning-based colour constancy~\cite{xie2023camera,kang2026graph,liu2026illuminant} improves benchmarks but needs large annotated corpora and domain retraining.
Multimodal ambient-sensor fusion~\cite{liu2026illuminant} and FAID-trained flash decomposition~\cite{flashambient,upadhyay2025low} suggest hybrid physical--learned pipelines.
We position our work as a \emph{physically grounded} flash/no-flash reflectance path with \emph{offline} target-based camera characterisation, evaluated on skin rather than general scenes.

\subsection{Flash/No-Flash Illuminant and Reflectance Estimation}

DiCarlo et al.~\cite{dicarlo2001illuminating} isolated flash by subtraction; Petschnigg et al.~\cite{petschnigg2004digital} formalised flash/no-flash pairs for enhancement.
Lu and Drew~\cite{lu2006practical} recovered ambient chromaticity without flash SPD or camera calibration at runtime.
Hui et al.~\cite{hui2016white,hui2018illuminant} applied separation to white balance and per-pixel illuminant spectra.
Maralan et al.~\cite{maralan2023computational} jointly estimated albedo, shading, and illuminant under a structured model.

Our primary reflectance estimate uses the geometric mean $\Rvec=\sqrt{\mathbf{A}\odot\mathbf{B}}$ after ECC alignment---a practical recipe related to but not identical with Lu's log-chroma locus~\cite{lu2006practical,evangelidis2008parametric}.
Camera RGB$\rightarrow$XYZ is supplied by an offline matrix fit on ColorChecker training captures (chart present at train time, absent at inference), following the spirit of Choi et al.'s skin-weighted target calibration~\cite{choi2017performance}.

\subsection{Skin Tone Measurement and Reference Scales}

Camera skin colour accuracy depends on illuminant characterisation, detector calibration, and ROI choice~\cite{harville2005consistent,choi2017performance,cook2025colorimetric}.
The Fitzpatrick scale~\cite{fitzpatrick1988} is ordinal and subjective; MST~\cite{cheng2024monastic} provides D65 Lab references suitable for device-independent comparison.
We use FitSkin contact-scanner cheek Lab as ground truth rather than MST assignment, but report in CIELAB under D65 for compatibility with MST-oriented literature.

%=============================================================================
\section{Methods}\label{sec:methods}
%=============================================================================

\subsection{Cohort and Ground Truth}

We study two participants imaged on iPhone ProRAW (DNG) with paired flash/no-flash captures:
\textbf{P1 (Emily)}, lighter skin (Caucasian-range FitSkin cheek $L^*\approx 69$); and
\textbf{P2 (Likitha / Liki)}, darker skin, \textbf{Indian ancestry} (FitSkin cheek $L^*\approx 52$).
The primary evaluation cohort (\emph{Pansor}, 16~June~2026) comprises six ColorChecker trials (three per participant) and six Sephora Bag trials excluded from the present chart-colorimetry analysis.
Ground truth is FitSkin multi-point \emph{cheek} CIELAB under D65, taken as the May~20 median cheek Lab per participant (\emph{cross-session} reference; iPhone capture ${\sim}4$ weeks later).
Agreement is CIEDE2000 ($\deee$) between pipeline cheek or mesh means and FitSkin cheek means.
Cross-session reference inflates absolute $\deee$ relative to same-session booth validation (Phase~4 median $\deee=3.50$, $N=5$) but preserves \emph{relative} comparison across pipelines.

Additional ColorChecker training sources (not used in every eval row):
\begin{itemize}[noitemsep]
  \item Phase-4 booth RAW DNGs (May 2025, six pairs);
  \item bundled chart\_cc JPEG pairs (six pairs, 8-bit export).
\end{itemize}

\subsection{RAW Ingestion}

DNGs were demosaiced to linear camera RGB (\texttt{rawpy}, half-resolution).
Pansor ProRAW required embedded \emph{camera white balance}; unity WB caused severe failure.
A per-channel 99.5th-percentile scale mapped linear RGB to $[0,1]$ for alignment; statistics used linear values before display encoding.

\subsection{Pipeline A: In-Scene ColorChecker (\texttt{run\_chart\_cc.py})}

When a 24-patch Macbeth ColorChecker is visible:
\begin{enumerate}[noitemsep]
  \item Detect MCC on the no-flash frame; apply gray-world white balance from neutral patches.
  \item Estimate RGB$\rightarrow$XYZ via least squares (lighter skin: $3\times3$) or affine $3\times4$ (darker skin).
  \item Extract mean CIELAB in a skin ROI under D65.
\end{enumerate}

\paragraph{Skin-tone auto policy.}
Preview cheek $L^*$ (gray-WB, uncorrected) routes capture settings:
\begin{align*}
  L^*_{\mathrm{probe}} \ge 42 &\Rightarrow \text{\textbf{light tier}: cheek ROI, } 3\times3 \text{ matrix},\\
  L^*_{\mathrm{probe}} < 42 &\Rightarrow \text{\textbf{dark tier}: face-mesh ROI, affine matrix}.
\end{align*}
On Pansor, P1 probe $L^*\approx 47$--$50$; P2 probe $L^*\approx 35$ (Table~\ref{tab:probe}).

\subsection{Pipeline B: Chart-Free Flash/No-Flash}

When no chart is in frame, reflectance is recovered from aligned flash/no-flash linear RGB:
\begin{equation}\label{eq:reflectance}
  \Rvec = \sqrt{\max(\mathbf{A}\odot\mathbf{B},\,\varepsilon)},
\end{equation}
where $\mathbf{A}$ is no-flash and $\mathbf{B}$ flash after ECC registration and \emph{skin-mask median-luma} exposure matching (flash scaled to no-flash on a MediaPipe face-mesh mask, not global scene scaling).
$\Rvec$ is mapped to CIELAB via offline $\Maff$ trained on separate ColorChecker sessions.

\paragraph{ROI rule (chart-free).}
\textbf{Face-mesh ROI only} for all skin tones in flash/no-flash mode.
Cheek-hull ROI on darker skin inflated $\deee$ to ${\sim}8$--$9$ (Section~\ref{sec:results-roi}).

\paragraph{Tone-routed offline matrices.}
Two bundles are maintained:
\begin{itemize}[noitemsep]
  \item \texttt{calibration/tier3\_by\_tone/light/} --- Emily training sessions + \texttt{issa\_median\_caucasian};
  \item \texttt{calibration/tier3\_by\_tone/dark/} --- Likitha Pansor ProRAW CC + \texttt{issa\_median\_south\_asian}.
\end{itemize}
At inference, \texttt{--skin-tone auto} selects the bundle by manifest participant or preview $L^*$.

\subsection{Offline Matrix Training}

Training uses flash/no-flash pairs \emph{with} visible ColorChecker (no chart required at deployment).
For each training trial, 24 MCC patches on the aligned no-flash frame provide camera-linear RGB; canonical D65 XYZ~\cite{BabelColorMCC} provides reference.
Patches are stacked across trials; optional ISSA median reflectance rows~\cite{issa2025} are synthesised through monochromator-derived camera spectral sensitivity.
Affine $\Maff\in\mathbb{R}^{4\times 3}$ is fit by weighted least squares with anchor down-weighting; Lu neutral-patch log-chroma sharpening~\cite{lu2006practical} is stored for optional use.

\paragraph{Per-row white balance.}
Pansor ProRAW training rows use \texttt{raw\_camera\_wb=yes}; booth RAW and JPEG rows use unity/camera-default as appropriate (Table~\ref{tab:training-manifest}).

\paragraph{Tone-specific training policy (empirical).}
\begin{itemize}[noitemsep]
  \item \textbf{Light bundle:} stack all nine Emily ColorChecker sources (3 Pansor + 3 booth + 3 JPEG).
  \item \textbf{Dark bundle:} Pansor ProRAW only (3 Likitha trials); booth/JPEG stacking raised inferred $L^*$ from ${\sim}50$ to ${\sim}64$ on ProRAW eval by mixing capture pipelines.
\end{itemize}

Training is reproduced by:
\begin{verbatim}
python3 scripts/train_skin_tone_bundles.py
\end{verbatim}

\subsection{Pipeline C: Hybrid Production Rule}

\begin{algorithmic}[1]
\If{ColorChecker detected in frame}
  \State Run Pipeline~A with \texttt{--skin-tone auto}
\ElsIf{flash/no-flash pair available}
  \State Probe skin-tone tier; load \texttt{light/} or \texttt{dark/} bundle
  \State Run Pipeline~B (mesh ROI, Eq.~\ref{eq:reflectance})
\EndIf
\end{algorithmic}

Optional Sephora-bag CAT02~\cite{fitskin2026pansor,cie2004cat,Bradford1984} applies on bag trials only; excluded from ColorChecker primary analysis here.

\subsection{Evaluation Metrics}

Primary endpoint: cheek or mesh mean $\deee$~\cite{sharma2005ciede2000} vs.\ FitSkin cheek Lab.
Secondary: signed $\Delta L^*$, $\Delta a^*$, $\Delta b^*$.
We report per-trial, per-participant, and cohort median/mean.

%=============================================================================
\section{Results}\label{sec:results}
%=============================================================================

\subsection{In-Scene ColorChecker: Fixed vs.\ Skin-Tone Auto}
\label{sec:results-chart}

Table~\ref{tab:chart-cc} compares fixed cheek ROI + single matrix fit (baseline) against skin-tone auto on six Pansor ColorChecker DNG trials.

\begin{table}[htbp]
\centering
\caption{In-scene ColorChecker on Pansor ProRAW ($N=6$): fixed cheek/$3\times3$ vs.\ \texttt{--skin-tone auto}. Ground truth: cross-session FitSkin May~20 cheek Lab.}
\label{tab:chart-cc}
\small
\begin{tabular}{l S[table-format=1.2] S[table-format=1.2] S[table-format=1.2] S[table-format=1.2]}
\toprule
{Trial} & {$\deee$ (fixed)} & {$\deee$ (auto)} & {Tier} & {$L^*_{\mathrm{probe}}$} \\
\midrule
P1\_CC\_T1 & 5.04 & 5.04 & light & 46.4 \\
P1\_CC\_T2 & 2.25 & 2.25 & light & 48.6 \\
P1\_CC\_T3 & 1.09 & 1.09 & light & 49.7 \\
P2\_CC\_T1 & 4.39 & \bfseries 0.96 & dark  & 34.8 \\
P2\_CC\_T2 & 5.95 & \bfseries 1.73 & dark  & 35.0 \\
P2\_CC\_T3 & 4.89 & \bfseries 0.93 & dark  & 34.6 \\
\midrule
{Median}   & 4.64 & \bfseries 1.41 & --- & --- \\
{Mean}     & 3.93 & \bfseries 2.13 & --- & --- \\
\bottomrule
\end{tabular}
\end{table}

\begin{table}[htbp]
\centering
\caption{Preview cheek $L^*$ and auto routing on Pansor ColorChecker trials.}
\label{tab:probe}
\begin{tabular}{lccS[table-format=1.1]}
\toprule
{Participant} & {Tier} & {ROI (auto)} & {$L^*_{\mathrm{probe}}$} \\
\midrule
P1 (Emily)   & light & cheek & 46.4--49.7 \\
P2 (Likitha) & dark  & mesh  & 34.6--35.0 \\
\bottomrule
\end{tabular}
\end{table}

Auto routing primarily benefits P2: mesh + affine reduces $\deee$ from ${\sim}4.4$--$5.9$ to ${\sim}0.9$--$1.7$.
P1 trials already used cheek ROI; residual error on P1\_CC\_T1 ($\deee=5.04$) reflects cross-session $L^*$ under-exposure on no-flash rather than ROI choice.

\subsection{Chart-Free Flash/No-Flash: Pooled vs.\ Tone-Routed}
\label{sec:results-roi}

Table~\ref{tab:fnf} summarises chart-free inference (mesh ROI, Eq.~\ref{eq:reflectance}) on the same six ColorChecker trials (chart \emph{not} used at inference; trials serve as labelled portrait captures).

\begin{table}[htbp]
\centering
\caption{Chart-free flash/no-flash ($N=6$, mesh ROI): single pooled offline matrix vs.\ tone-routed \texttt{light/}+\texttt{dark/} bundles. Final training: light = 9 Emily sessions; dark = 3 Pansor ProRAW + ISSA South Asian.}
\label{tab:fnf}
\small
\begin{tabular}{l S[table-format=1.2] S[table-format=1.2] S[table-format=1.2]}
\toprule
{Condition} & {P1 median $\deee$} & {P2 median $\deee$} & {All median $\deee$} \\
\midrule
Pooled \texttt{tier3\_affine}     & 2.03 & 5.26 & 3.65 \\
Tone-routed (v1: 3-trial bundles) & 5.77 & \bfseries 2.16 & 4.11 \\
\textbf{Tone-routed (final)}       & \bfseries 1.77 & \bfseries 2.16 & \bfseries 2.16 \\
Cheek ROI on P2 (FNF, pooled)     & ---  & ${\sim}8.8$ & --- \\
\midrule
Chart CC + auto (reference)       & 1.67 & 1.35 & 1.41 \\
\bottomrule
\end{tabular}
\end{table}

\begin{table}[htbp]
\centering
\caption{Per-trial chart-free $\deee$ with final tone-routed bundles (mesh ROI).}
\label{tab:fnf-per-trial}
\begin{tabular}{l S[table-format=1.2] S[table-format=1.2] S[table-format=1.2]}
\toprule
{Trial} & {Pooled $\deee$} & {Tone-routed $\deee$} & {$L^*_{\mathrm{pipe}}$ (tone)} \\
\midrule
P1\_CC\_T1 & 2.31 & 2.39 & 66.0 \\
P1\_CC\_T2 & 2.03 & 1.77 & 66.8 \\
P1\_CC\_T3 & 1.93 & 1.69 & 67.0 \\
P2\_CC\_T1 & 5.26 & 2.75 & 50.2 \\
P2\_CC\_T2 & 5.69 & 1.94 & 50.9 \\
P2\_CC\_T3 & 5.00 & 2.45 & 49.9 \\
\bottomrule
\end{tabular}
\end{table}

Tone routing corrected P2 $L^*$ from ${\sim}58$ (pooled) to ${\sim}50$ (FitSkin ${\sim}52$).
A single pooled matrix traded off ${\sim}5\,\deee$ on darker skin against ${\sim}2\,\deee$ on lighter skin.

\subsection{Training-Data Ablation (Dark Bundle)}

Table~\ref{tab:dark-train} shows P2 median chart-free $\deee$ when varying dark-bundle training sources (mesh ROI eval on Pansor).

\begin{table}[htbp]
\centering
\caption{Dark-bundle training ablation: P2 median chart-free $\deee$ on Pansor ($N=3$).}
\label{tab:dark-train}
\begin{tabular}{l S[table-format=1.0] S[table-format=1.2]}
\toprule
{Training sources} & {$n_{\mathrm{trials}}$} & {P2 median $\deee$} \\
\midrule
Pansor ProRAW only + ISSA           & 3 & \bfseries 2.16 \\
Pansor + booth RAW + ISSA           & 6 & 4.40 \\
Pansor + booth RAW + JPEG + ISSA    & 9 & 10.3 \\
\bottomrule
\end{tabular}
\end{table}

Booth RAW (unity WB, May session) and 8-bit JPEG exports are valuable for lighter-skin stacking but act as \emph{negative transfer} when naively merged into a darker-skin ProRAW matrix.

\begin{table}[htbp]
\centering
\caption{Light-bundle training manifest (nine trials).}
\label{tab:training-manifest}
\small
\begin{tabular}{llcc}
\toprule
{subject\_id} & {source} & {camera WB} & {participant} \\
\midrule
P1\_CC\_T1--T3     & Pansor ProRAW & yes & Emily \\
P1\_T1--T3\_booth  & Booth RAW     & no  & Emily \\
P1\_JPEG\_T1--T3   & Chart JPEG    & no  & Emily \\
\bottomrule
\end{tabular}
\end{table}

\subsection{ISSA Spectral Prior: South Asian vs.\ African}

Likitha (P2) is Indian; we replaced \texttt{issa\_median\_african} with \texttt{issa\_median\_south\_asian} in the dark bundle and SCR-AWB defaults.
MCC patches dominate the matrix fit; prior choice shifts P2 chart-free $\deee$ modestly (Table~\ref{tab:issa}).

\begin{table}[htbp]
\centering
\caption{ISSA prior ablation on dark bundle (Pansor P2, chart-free).}
\label{tab:issa}
\begin{tabular}{l S[table-format=1.2] S[table-format=1.2] S[table-format=1.2] S[table-format=1.2]}
\toprule
{ISSA prior} & {T1 $\deee$} & {T2 $\deee$} & {T3 $\deee$} & {Median} \\
\midrule
\texttt{issa\_median\_african}      & 2.75 & 1.94 & 2.45 & 2.45 \\
\textbf{\texttt{issa\_median\_south\_asian}} & 2.67 & 1.89 & 2.33 & \bfseries 2.33 \\
\bottomrule
\end{tabular}
\end{table}

\paragraph{SCR-AWB baseline.}
Spectral-reflectance auto white balance~\cite{zhou2025scr} with ISSA prior alone achieved median $\deee\approx 13.8$ on Pansor and is not sufficient as a standalone skin colorimeter.

%=============================================================================
\section{Discussion}\label{sec:discussion}
%=============================================================================

\subsection{Relation to Prior Work}

Our results align with Choi et al.~\cite{choi2017performance}: skin-colour accuracy improves when calibration and ROI choices target the skin region rather than treating all reflectances equally.
Extending this principle across \emph{tone}, we find that darker skin benefits from mesh sampling and affine RGB$\rightarrow$XYZ, consistent with spatial non-uniformity and specular/sub-surface effects at cheek hull boundaries~\cite{harville2005consistent}.

The flash/no-flash geometric mean (Eq.~\ref{eq:reflectance}) operationalises the active illumination idea of DiCarlo~\cite{dicarlo2001illuminating} and Lu~\cite{lu2006practical} without requiring runtime flash SPD estimation, at the cost of a separate offline camera characterisation---an instance of \emph{train with chart, deploy without chart}, distinct from classical in-scene ColorChecker correction~\cite{wakholi2026systematic} and from fully targetless grey-world methods~\cite{gijsenij2011}.

Tone-routed offline matrices address a failure mode not visible in single-participant calibrations: pooled fits compromise darker skin $L^*$ by ${\sim}6$ units in our cohort.
This resonates with fairness critiques of Fitzpatrick-thresholded systems~\cite{cook2025colorimetric} and motivates device-independent scales (MST~\cite{cheng2024monastic}) \emph{and} capture-pipeline stratification.

Separating \textbf{lightness tier} (preview $L^*$ routing) from \textbf{ancestry prior} (ISSA South Asian vs.\ Caucasian) clarifies design for Indian and other South Asian users: tier controls ROI/matrix class; prior anchors spectral training rows~\cite{issa2025,zhou2025scr,wang2025perform}.

\subsection{Commercial Deployment}

Table~\ref{tab:commercial} summarises realistic product tiers without a ColorChecker in every capture.

\begin{table}[htbp]
\centering
\caption{Commercial feasibility without per-capture ColorChecker.}
\label{tab:commercial}
\small
\begin{tabular}{p{0.28\linewidth} c p{0.18\linewidth} p{0.32\linewidth}}
\toprule
{Use case} & {Feasible?} & {Expected $\deee$} & {Requirements} \\
\midrule
Shade matching / relative tone & Yes & 2--4 & Fixed phone model, FNF RAW, mesh ROI, tone-routed matrix \\
Longitudinal skincare trends   & Yes & Relative & Same-user baseline, controlled capture protocol \\
Absolute Lab vs.\ FitSkin      & Hard & $>$2 without CC & In-scene reference or same-session contact scan \\
\bottomrule
\end{tabular}
\end{table}

A practical architecture mirrors Pipeline~C (Section~\ref{sec:methods}): factory/calibration sessions with ColorChecker; daily capture via tone-routed FNF; optional bag or gray-card refinement~\cite{fitskin2026pansor}.
Product copy should emphasise \emph{shade index} or change-from-baseline rather than absolute CIELAB without reference.

Conditions that \emph{cannot} be ignored:
\begin{itemize}[noitemsep]
  \item one global matrix across tones and phone models;
  \item social-filter JPEG capture;
  \item cheek ROI on chart-free darker skin;
  \item ancestry-mismatched ISSA priors (e.g.\ African prior for Indian skin).
\end{itemize}

\subsection{Limitations}

\begin{itemize}[noitemsep]
  \item $N=2$ participants; cross-session FitSkin reference inflates absolute $\deee$.
  \item Single iPhone model / ProRAW processing path; per-device matrices remain necessary.
  \item Dark-bundle training restricted to three Pansor trials until more homogenous ProRAW sessions are collected.
  \item Sephora-bag and Lu-booth paths excluded from ColorChecker primary analysis; bag CAT02 helps P2 bag trials but not P1~\cite{fitskin2026pansor}.
\end{itemize}

\subsection{Future Work}

Collect additional Indian/darker ProRAW ColorChecker sessions; weighted multi-source training; same-session FitSkin validation; MST scale assignment from corrected Lab; per-device matrix shipping; evaluation on bag-only (chart-free) trials.

%=============================================================================
\section{Conclusion}\label{sec:conclusion}
%=============================================================================

Non-contact iPhone skin colorimetry is feasible at $\deee\approx 2$ chart-free and $\deee\approx 1.4$ with in-scene ColorChecker when calibration is \textbf{tone-routed} and \textbf{ancestry-aware}.
Likitha's Indian skin required South Asian ISSA priors and a dark-tier mesh/affine policy; Emily's lighter skin benefited from stacking nine ColorChecker training sessions.
A hybrid deployment---chart when visible, flash/no-flash otherwise---balances FitSkin-grade accuracy with consumer capture constraints.
Absolute contact-scanner agreement without any in-scene reference remains hard; shade matching and longitudinal monitoring are the near-term commercial targets.

%=============================================================================
\section*{Reproducibility}
%=============================================================================

\noindent
Repository: \url{https://github.com/RooneyEmily/Fitskin}.
Key commands:
\begin{verbatim}
python3 scripts/build_pansor_manifest.py --data-root "$PANSOR_DATA_ROOT"
python3 scripts/train_skin_tone_bundles.py
python3 run_chart_cc.py --input-mode dng \
  --manifest data/pansor/manifest_pansor_fitskin.csv --cc-only \
  --skin-tone auto --out-dir chart_cc_output/pansor_dng_auto_tone
python3 flash_no_flash_skin_lab.py --input-mode dng \
  --manifest data/pansor/manifest_pansor_fitskin.csv \
  --iphone-calibration-tone-root calibration/tier3_by_tone \
  --skin-tone auto --raw-camera-wb --exposure-scale-skin-mask \
  --bag-cat02 off --production
\end{verbatim}

%=============================================================================
% Use your Overleaf .bib file. Append entries from references_supplement.bib
% for keys not yet in your library (wu1999, issa2025, zhou2025scr, fitskin2026pansor).
% Compile: pdflatex -> bibtex -> pdflatex -> pdflatex
%=============================================================================
\bibliographystyle{plain}
\bibliography{references,references_supplement}

\end{document}
